\documentclass[11pt]{article}

\usepackage[margin=1.5in]{geometry}
\usepackage{lecnotes}
% \usepackage{mpass}
% \lstset{style=verb,language=mpass}
\input{lmacros}

\newcommand{\course}{15-316: Software Foundations of Security \& Privacy}
\newcommand{\lecturer}{Matt Fredrikson}
\newcommand{\lecdate}{March 19, 2026}
\newcommand{\lecnum}{17}
\newcommand{\lectitle}{Proof Search in Authorization Logic}
\newcommand{\courseurl}{https://15316-cmu.github.io/2026/}

\begin{document}

\maketitle

\section{Introduction}

In Boolean logic, the sequent calculus was a good basis for proof search (at
least on the small scale) because all rules were sound and invertible.  This is
no longer the case for intuitionistic logic, whether extended with affirmation
or not.  In fact, Boolean propositional logic is just about the only logic where
we can arrange for all rules to be invertible.  Our task then is to find out
which rules \emph{are} invertible and which are not.  This gives rise to a
classification of logical connectives based on the invertibility of their right
and left rules.  This can be done generally for logics that admit a sequent
calculus formulation.  It was first discovered for \emph{linear logic}
\citep{Girard87tcs} by \citet{Andreoli92jlc} but since been applied to many
other ones, including those with a lax modality \citep{Liang09tcs,Watkins02tr}.

Our first order of business, then, is to identify invertibility properties,
followed by proof strategies based on them.

\section{Inversion}

Since it is a pervasively used connective, we start with implication.
\begin{rules}
  \begin{array}[t]{c}
    \infer[{\arrow}R]
    {\Gamma \vdash P \arrow Q}
    {\Gamma, P \vdash Q}
    \\[1em]
    \mbox{\color{red}invertible}
  \end{array}
  \hspace{3em}
  \begin{array}[t]{c}
    \infer[{\arrow}L]
    {\Gamma, P \arrow Q \vdash \delta}
    {\Gamma \vdash P
    & \Gamma, Q \vdash \delta}
    \\[1em]
    \mbox{\color{red}not invertible}
  \end{array}
\end{rules}
It turns out the right rule is invertible, while the left rule is not.  To see
that the left rule is \emph{not} invertible, we only need a counterexample.
Consider
\[
  p \arrow q, q \arrow r \vdash p \arrow r
\]
Trying to apply the left rule for implication to the first antecedent
results in a sequent that cannot be derived.
\[
  \infer[{\arrow}L]
  {p \arrow q, q \arrow r \vdash p \arrow r}
  {\infer[{\arrow}L]
    {q \arrow r \vdash p}
    {\deduce[\ms{XXX}]{\cdot \vdash q}{}
      & \deduce[\ms{XXX}]{r \vdash p}{}}
    & \deduce[\ddd]{q, q \arrow r \vdash p \arrow r}{}}
\]
We suspect that the right rule for implication is invertible, but how do we
prove it?  Previously, we used a \emph{semantic argument}, using the validity of
a sequent.  Here, such a direct argument is not available---we would have to
introduce a semantics which is not straightforward.  But there is also a
syntactic technique, using the \emph{admissibility of cut}.  So we make a brief
detour to introduce it.

When presented as a rule
\begin{rules}
  \infer[\ms{cut}]
  {\Gamma \vdash \delta}
  {\Gamma \vdash P & \Gamma, P \vdash \delta}
\end{rules}
its meaning is clear: if we can prove $P$ we are allowed to assume it in further
reasoning.  For bottom-up proof search it is somewhat problematic because we have
to determine a useful $P$, which could be \emph{any} formula.  So we use it only
under controlled circumstances.

\citet{Gentzen35} showed that the rule of cut is redundant for both classical
(Boolean) and intuitionistic logic.  What do we mean by ``redundant''?  More
technically, we call a rule \emph{admissible} if it is both sound and whenever
the premises can be derived, so can the conclusion \emph{without using the
  rule}.  We write an admissible rule with a dashed line:
\begin{rules}
  \infer-[\ms{cut}]
  {\Gamma \vdash \delta}
  {\Gamma \vdash P & \Gamma, P \vdash \delta}
\end{rules}
We won't go into detail how to show that cut is admissible.  This is covered in
many articles (including Gentzen's original one) and also in \emph{15-317
  Constructive Logic} in somewhat more modern notation (see
\href{https://www.cs.cmu.edu/~fp/courses/15317-s23/lectures/08-cutelim.pdf}{Lecture
  8}).

Here we concentrate on how to use the admissibility of cut to obtain other
properties.  First, it is convenient to have an alternative version that is also
admissible and more suitable to top-down reasoning.
\begin{rules}
  \infer-[\ms{cut}']
  {\Gamma_1, \Gamma_2 \vdash \delta}
  {\Gamma_1 \vdash P & \Gamma_2, P \vdash \delta}
\end{rules}
We want to prove that ${\arrow}R$ is invertible.  That is, we want to show that
\[
  \infer-[{\arrow}R^{-1}]
  {\Gamma, P \vdash Q}
  {\Gamma \vdash P \arrow Q}
\]
is admissible.  Step-by-step: first, we have a succedent $P \arrow Q$ in the
premise, so we should try to cut it with a sequent with antecedent $P \arrow Q$
with some $\Gamma'$ and $\delta$.
\[
  \infer-[\ms{cut}']
  {\Gamma, P \vdash Q}
  {\Gamma \vdash P \arrow Q
    & \deduce[\ddd]{\Gamma', P \arrow Q \vdash \delta}{}}
\]
What could we choose for $\Gamma'$ and $\delta$?  It is pretty obvious from the
conclusion: $\Gamma' = P$ and $\delta = Q$.  We can then complete the derivation
of the second premise.
\[
  \infer-[\ms{cut}']
  {\Gamma, P \vdash Q}
  {\Gamma \vdash P \arrow Q
    & \infer[{\arrow}L]
    {P, P \arrow Q \vdash Q}
    {\infer[\ms{id}]{P \vdash P}{}
      & \infer[\ms{id}]{Q \vdash Q}{}}}
\]
We see in the second premise everything is proved, so this derivation shows that
${\arrow}R^{-1}$ is admissible.

We can prove other rules invertible as well, with clever uses of $\ms{cut}'$.
On our fragment (excluding affirmation for now), the invertible rules are
${\land}R$, ${\land}L$, ${\lor}L$, ${\forall}R$.  The noninvertible rules are
${\arrow}L$, ${\lor}R_1$, ${\lor}R_2$, and ${\forall}L$.

A first simple strategy emerges: apply all the invertible rules in some
arbitrary, unspecified order until we reach a sequent where either identity
applies, or we have to make a choice between noninvertible rules.  This choice
typically then requires backtracking, in case we make a wrong choice.

We can express such a strategy with three judgment forms, two that force only
invertible rules to be used on the left or right, and one that requires a
choice.  Since it is not our ultimate destination, we elide this here and refer
the interest reader to
\href{https://www.cs.cmu.edu/~fp/courses/15317-s23/lectures/15-inversion.pdf}{Lecture
  15} of the course on
\href{https://www.cs.cmu.edu/~fp/courses/15317-s23}{constructive logic}.

\section{Inversion for Affirmation}

An empirical observation is that if the right rule for a connective is
invertible then the left rule is not and vice versa.  This is apparently
violated by conjunction which is invertible on both sides.  This is because
there are actually two forms of conjunction hiding under the symbol ``$\land$''
that are indistinguishable with respect to provability.

We call connectives whose right rule is invertible \emph{negative connectives}
and the ones whole left rule is invertible are \emph{positive connectives}.  The
inversion phase of proof search then applies negative right rules and positive
left rules.

But what about affirmation?  There is something rather strange going on because
in the right rule we go from $(A \says P)\; \ms{true}$ to $A \aff P$ and in the
left rule we jump directly from $A \says P$ to $P$.  It turns out that the
affirmation modality signifies a change in polarity.  One way to state that is
that $A \says P$ is negative, but the judgment underneath, $A \aff P$ is
positive.  That means, the right rule is invertible, but the rule of affirmation
is not.
\begin{rules}
  \begin{array}[t]{c}
    \infer[{\mb{says}}R]
    {\Gamma \vdash (A \says P)\; \ms{true}}
    {\Gamma \vdash A \aff P}
    \\[1em]
    \mbox{\color{red}invertible}
  \end{array}
  \hspace{3em}
  \begin{array}[t]{c}
    \infer[\mb{aff}]
    {\Gamma \vdash A \aff P}
    {\Gamma \vdash P\; \ms{true}}
    \\[1em]
    \mbox{\color{red}not invertible}
  \end{array}
\end{rules}
As for the left rules, because $A \says P$ is negative, it's left rules is not
invertible.  However, if the succedent has the form $A \aff Q$ we can always
apply the rule since $P$ is a stronger assumption than $A \says P$.  So we put
``invertible'' in quotation: we can apply the rule when possible, but we cannot
always apply it when $A \says P$ is an antecedent.
\begin{rules}
  \begin{array}[t]{c}
    \infer[\mb{says}L]
    {\Gamma, A \says P \vdash A \aff Q}
    {\Gamma, P \vdash A \aff Q}
    \\[1em]
    \mbox{\color{red}``invertible''}
  \end{array}
\end{rules}
Now we also see why $A \aff P$ does not appear as a judgment among the
antecedents: it is positive, and therefore the corresponding rule can be applied
immediately and the stepping stone be omitted.
\begin{rules}
  \infer[\mb{says}L]
  {\Gamma, A \says P \vdash A \aff Q}
  {\infer[]
    {\Gamma, \mbox{``$A \aff P$''} \vdash A \aff Q}
    {\Gamma, P \vdash A \aff Q}}
\end{rules}

\section{Focusing}

If we look at the example from the last lecture we see that the only invertible
step is the right rule for $\mb{says}$, followed by stripping the
``$\mi{admin} \says \null$'' prefix from $(1)$, $(2)$, and $(3)$ using
$\mb{says}L$.
\begin{tabbing}
  \quad $(1) : \mi{admin} \says (\forall A.\, \forall R.\, \ms{owns}(A, R) \arrow \ms{mayOpen}(A, R))$, \\
  \quad $(2) : \mi{admin} \says (\forall A.\, \forall B.\, \forall R.\, \ms{owns}(A, R)
  \land \mi{mfredrik} \says \ms{studentOf}(B, A) \arrow \ms{mayOpen}(B, R))$, \\[1ex]
  \quad $(3) : \mi{admin} \says \ms{owns}(\mi{mfredrik}, \mi{cic2126})$, \\
  \quad $(4) : \mi{mfredrik} \says \ms{studentOf}(\mi{saranyav}, \mi{mfredrik})$ \\[1ex]
  \quad $\vdash$ \\[1ex]
  \quad $\mi{admin} \says \ms{mayOpen}(\mi{saranyav}, \mi{cic2126})$
\end{tabbing}

\clearpage
\noindent The resulting sequent below has no invertible rule we can blindly apply.
\begin{tabbing}
  \quad $(1)' : \forall A.\, \forall R.\, \ms{owns}(A, R) \arrow \ms{mayOpen}(A, R)$, \\
  \quad $(2)' : \forall A.\, \forall B.\, \forall R.\, \ms{owns}(A, R)
  \land \mi{mfredrik} \says \ms{studentOf}(B, A) \arrow \ms{mayOpen}(B, R)$, \\[1ex]
  \quad $(3)' : \ms{owns}(\mi{mfredrik}, \mi{cic2126})$, \\
  \quad $(4) : \mi{mfredrik} \says \ms{studentOf}(\mi{saranyav}, \mi{mfredrik})$ \\[1ex]
  \quad $\vdash$ \\[1ex]
  \quad $\mi{admin} \aff \ms{mayOpen}(\mi{saranyav}, \mi{cic2126})$
\end{tabbing}
For example, we could instantiate the quantifier $\forall A$ in $(1)'$.  But
still, everything remains negative and we could instantiate either the
$\forall R$ we have uncovered or the $\forall A$ quantifier in $(2)$.  We see
that at every step we have to choose between a number of alternatives.  This can
lead to a lot of backtracking if the choices are incorrect.

\emph{Focusing} \citep{Andreoli92jlc} is the idea that we can pick a negative
antecedent or a positive succedent and continue to apply noninvertible rules to
this particular formula until we reach either an atomic formula or the polarity
switches.  For example, if we guess $(2)'$ we would apply ${\forall}L$ three
times, followed by ${\arrow}L$.  Let's treat conjunction as positive, which we
continue with that as well and succeed when we find the needed
assumptions.\footnote{Actually, not quite.  The affirmation ``$\mi{mfredrik} \says$''
  forces us to stop and make another explicit choice.}  All that reasoning adds
$\ms{mayOpen}(\mi{saranyav}, \mi{cic2126})$ to the assumptions and we have to make
another choice.  At this point we want to choose the rule of affirmation
followed by the identity to complete the proof.

We don't write out the sequent calculus for focusing in its most general form,
but you may refer to
\href{https://www.cs.cmu.edu/~fp/courses/15317-s23/lectures/17-focusing.pdf}{Lecture
  17} of the constructive logic course notes for details.  We write
$\Gamma, [P] \vdash \delta$ for a formula $P$ in left focus and
$\Gamma \vdash [Q]$ for a formula in right focus.  In this kind of sequent
just one formula can be in focus, and rules can be applied only to the formula
in focus.   The rules with no focus are the invertible rules.

This gives us the following rules.  We work with following polarities below.
Except for atoms and conjunction (which we choose to be positive), they are
determined by the inversion properties of the connectives.
\[
  \begin{array}{llcl}
    \mbox{Negative} & P^-, Q^- & ::= & P \arrow Q \mid \forall x.\, P(x) \mid A \says P \\
    \mbox{Positive} & P^+, Q^+ & ::= & p \mid P \land Q \mid P \lor Q
  \end{array}
\]
First the rules to \emph{initiate} a phase of focusing.
\begin{rules}
  \infer[\mb{focus}R]
  {\Gamma \vdash Q^+}
  {\Gamma \vdash [Q^+]
    % & \mbox{($Q^+$ positive)}
  }
  \hspace{3em}
  \infer[\mb{focus}L]
  {\Gamma \vdash \delta}
  {P^- \in \Gamma & \Gamma, [P^-] \vdash \delta
    % & \mbox{($P^-$ negative)}
  }
\end{rules}
In the left rule we focus on a \emph{copy} of $P^-$ because we may need this
formula again.  In the examples we sometimes drop the extra copy if we
anticipate (or know) it will not be needed again.

The rule for affirmation is an additional transition rule.  Next, the logical
rules.  For brevity, we mostly omit ``$\ms{true}$'' in the succedent if it is
just a formula.  The judgment $\Gamma \vdash \delta$,
$\Gamma, [P] \vdash \delta$ and $\Gamma \vdash [Q]$ are mutually exclusive, that
is, there may be no focused formula in $\Gamma$ or $\delta$.
\begin{rules}
  \infer[\ms{id}]
  {\Gamma, p \vdash [p]}
  {}
  \\[1em] % \hline\\[1ex]
  \infer[{\arrow}R]
  {\Gamma \vdash P \arrow Q}
  {\Gamma, P \vdash Q}
  \hspace{3em}
  \infer[{\arrow}L]
  {\Gamma, [P \arrow Q] \vdash \delta}
  {\Gamma \vdash [P]
    & \Gamma, [Q] \vdash \delta}
  \\[1em]
  \infer[{\land}R]
  {\Gamma \vdash [P \land Q]}
  {\Gamma \vdash [P]
    & \Gamma \vdash [Q]}
  \hspace{3em}
  \infer[{\land}L]
  {\Gamma, P \land Q \vdash \delta}
  {\Gamma, P, Q \vdash \delta}
  \\[1em]
  \infer[{\forall}R^y]
  {\Gamma \vdash \forall x.\, P(x)}
  {\Gamma \vdash P(y) & \mbox{$y \not\in \Gamma, P(x)$}}
  \hspace{3em}
  \infer[{\forall}L]
  {\Gamma, [\forall x.\, P(x)] \vdash \delta}
  {\Gamma, [P(c)] \vdash \delta}
  \\[1em]
  \infer[{\lor}R_1]
  {\Gamma \vdash [P \lor Q]}
  {\Gamma \vdash [P]}
  \hspace{1em}
  \infer[{\lor}R_2]
  {\Gamma \vdash [P \lor Q]}
  {\Gamma \vdash [Q]}
  \hspace{3em}
  \infer[{\lor}L]
  {\Gamma, P \lor Q \vdash \delta}
  {\Gamma, P \vdash \delta
    & \Gamma, Q \vdash \delta}
\end{rules}
Next, the rules to complete a focusing phase (still postponing affirmations).
We say that we \emph{blur the focus}.
\begin{rules}
  \infer[\ms{blur}L]
  {\Gamma, [P^+] \vdash \delta}
  {\Gamma, P^+ \vdash \delta
    % & \mbox{($P^+$ positive)}
  }
  \hspace{3em}
  \infer[\ms{blur}R]
  {\Gamma \vdash [Q^-]}
  {\Gamma \vdash Q^-
    % & \mbox{($Q^-$ negative)}
  }
\end{rules}
Finally, the rules for affirmation.  They incorporate some phase transitions
because of the polarity shift intrinsic to the modality.
\begin{rules}
  \infer[{\mb{says}}R]
  {\Gamma \vdash (A \says P)\; \ms{true}}
  {\Gamma \vdash A \aff P}
  \hspace{3em}
  \infer[{\mb{says}}L]
  {\Gamma, [A \says P] \vdash A \aff Q}
  {\Gamma, P \vdash A \aff Q}
  \\[1em]
  \infer[\mb{aff}]
  {\Gamma \vdash A \aff P}
  {\Gamma \vdash [P]\; \ms{true}}
\end{rules}

Our motivating example is too lengthy to write out formally with these rules for
focusing, but we can try a small example and see how focusing reduces
nondeterminism.  Consider
\[
  p \arrow q, p \arrow (q \arrow r) \vdash p \arrow r
\]
After the obligatory right inversion, we arrive at
\[
  p, p \arrow q, p \arrow (q \arrow r) \vdash r
\]
In this situation, we could in principle try to focus
on $p \arrow q$, on $p \arrow (q \arrow r)$ or $r$.
Let's try right focus first.  We immediately fail because
$r$ is not among the antecedents.
\[
  \infer[\mb{focus}R]
  {p, p \arrow q, p \arrow (q \arrow r) \vdash r}
  {\deduce[\ms{XXX}]{p, p \arrow q, p \arrow (q \arrow r) \vdash [r]}{}}
\]
Next we try to focus on $p \arrow (q \arrow r)$.
\[
  \infer[\mb{focus}L]
  {p, p \arrow q, p \arrow (q \arrow r) \vdash r}
  {\infer[{\arrow}L]
    {p, p \arrow q, [p \arrow (q \arrow r)] \vdash r}
    {\infer[\ms{id}]
      {p, p \arrow q \vdash [p]}
      {}
      & \infer[{\arrow}L]
      {p, p \arrow q, [q \arrow r] \vdash r}
      {\deduce[\ms{XXX}]{p, p \arrow q \vdash [q]}{}
        & \deduce[\ddd]{p, p\arrow q, [r] \vdash r}{}}}}
\]
Note that after we decided which antecedent to focus on, everything was
determined.  We fail, because $q$ is not among the antecedents.

We cannot focus on $p$ on the left because it is positive, but we
can try one final option, namely to focus on $p \arrow q$.
\[
  \infer[\mb{focus}L]
  {p, p \arrow q, p \arrow (q \arrow r) \vdash r}
  {\infer[{\arrow}L]
    {p, [p \arrow q], p \arrow (q \arrow r) \vdash r}
    {\infer[\ms{id}]
      {p, p \arrow (q \arrow r) \vdash [p]}
      {}
      & \infer[\mb{blur}L]
      {p, [q], p \arrow (q \arrow r) \vdash r}
      {\deduce[\ddd]{p, q, p \arrow (q \arrow r) \vdash r}{}}}}
\]
Now it is possible to focus on $p \arrow (q \arrow r)$ because both $p$ and $q$
are among the antecedents.  This will add $r$ to the antecedents and we can
finally successfully focus on the succedent.

Remarkably, with focusing there is just a single proof (assuming we don't copy
the formula in focus).

Using inversion and focusing is sound and complete with respect
to the sequent calculus.  Soundness is important

\begin{theorem}[Soundness and Completeness of Inversion and Focusing]
  $\Gamma \vdash \delta$ in the sequent calculus if and only if
  $\Gamma \vdash \delta$ in the calculus with inversion and focusing.
\end{theorem}
\begin{proof}
  Soundness is straightforward, since we only restrict the application of
  certain rules.

  The proof of completeness is quite complex, and not just because of
  affirmations.  See \citet{Liang09tcs} for a blueprint that can be adapted to
  this authorization logic.
\end{proof}

In the architecture of proof-carrying authorization we have to contend with the
fact that more complex policies engender more difficult theorem proving
problems.  In the next lecture we will identify a fragment of authorization
logic that represents a reasonable compromise between expressiveness and
difficulty of proving.  It is inspired by Horn clauses, but goes beyond it due
the presence of affirmations.

\section{Practice Problems}

These problems are for self-study and are not to be handed in.

\subsection{Inversion and Focusing [4 parts]}

Consider the following sequent:
\[
  p \arrow (q \land r),\; q \arrow s \;\vdash\; p \arrow s
\]

\begin{task}\rm
  \label{task:inv1}
  What are the invertible rules that apply to this sequent?  Apply them to
  obtain a sequent where no more invertible rules apply.
\end{task}

\begin{task}\rm
  \label{task:inv2}
  Classify each formula in the resulting sequent as positive or negative.
  Which formulas can you focus on?
\end{task}

\begin{task}\rm
  \label{task:inv3}
  Try focusing on $p \arrow (q \land r)$.  Write out the focused derivation.
  Does it succeed?  If it gets stuck, explain where and why.
\end{task}

\begin{task}\rm
  \label{task:inv4}
  Find a focusing strategy that completes the proof.  How many choice points
  are there in your derivation?
\end{task}

\subsection{Focusing with Affirmations [3 parts]}

Consider the following sequent, where $\mi{mfredrik}$, $\mi{admin}$, and
$\mi{saranyav}$ are principals:
\[
  \mi{admin} \says (\mi{mfredrik} \says \ms{approved}(\mi{saranyav}) \arrow \ms{access}(\mi{saranyav}))
  \;\vdash\;
  \mi{admin} \says (\mi{mfredrik} \says \ms{approved}(\mi{saranyav}) \arrow \ms{access}(\mi{saranyav}))
\]

\begin{task}\rm
  \label{task:aff1}
  This sequent is trivially provable by identity.  Now consider a more
  interesting variant:
  \begin{tabbing}
    \quad $(1) : \mi{admin} \says (\forall B.\, \mi{mfredrik} \says \ms{approved}(B) \arrow \ms{access}(B))$, \\
    \quad $(2) : \mi{mfredrik} \says \ms{approved}(\mi{saranyav})$ \\[1ex]
    \quad $\vdash$ \\[1ex]
    \quad $\mi{admin} \says \ms{access}(\mi{saranyav})$
  \end{tabbing}
  Apply all invertible rules.  What is the resulting sequent?  Pay attention
  to which antecedents can be unwrapped via $\mb{says}L$ and which cannot.
\end{task}

\begin{task}\rm
  \label{task:aff2}
  Starting from the sequent you obtained in the previous task, identify the
  formula to focus on.  Write out the focused derivation, showing each step.
  Where does the polarity shift from focusing to inversion?
\end{task}

\begin{task}\rm
  \label{task:aff3}
  Explain why the affirmation $(2)$ from $\mi{mfredrik}$ does not need to be
  unwrapped via $\mb{says}L$ in order for the proof to go through, even though
  the policy in $(1)$ refers to $\mi{mfredrik} \says \ms{approved}(B)$.
\end{task}

\bibliographystyle{plainnat}
\bibliography{bibliography}

\end{document}
